\newpage
\begin{center}
\normalfont\Large\bfseries MPC exercise: week six
\end{center}
\setcounter{section}{0}
\section{Off-set free MPC}
For an LTI system
\begin{equation*}
    x_{k+1} = Ax_k + Bu_k,
\end{equation*}
the standard MPC formulation is
\begin{align*}
        \min_u \quad &\sum_{k = 0}^{N-1} x_k^{\top}Qx_k + u_k^{\top}Ru_k + x_N^{\top}Px_N,\\
    \text{s.t.}\quad &x_{k+1} =Ax_k + Bu_k, \quad\forall k \in [N],\\
    &Fx_k\leq f, \quad\forall k \in [N-1],\\
    &Gu_k\leq g, \quad\forall k \in [N-1],\\
    &Hx_N\leq h,
\end{align*}
where $P$ is the solution to the discrete-time algebraic Riccati equation and the set $\{x|Hx\leq h\}$ is the corresponding terminal set (maximum invariant set given the optimal LQR controller). This formulation can regulate the system towards the origin with stability and recursive feasible guarantees.
\begin{enumerate}
    \item \label{ex6:1} We want to reformulate the above MPC formulation such that it can track a known and feasible reference $x_{\text{ref}}$. Assume we can measure the full state, write down the steady state condition for $u_s$. \vspace{5\baselineskip}

    \item Let $\Delta x_k = x_k - x_s$, $\Delta u_k = u_k-u_s$, and assume the terminal set under delta formulation is $H'\Delta x_N\leq h'$, write down the delta formulation and explain how to recover the real control inputs from the delta formulation.\vspace{5\baselineskip}
    \item \label{ex6:1} An unknown constant disturbance $d$ perturbs the LTI system, which makes the system $x_{k+1} = Ax_k + Bu_k + d$. Assume your state estimator returns a disturbance estimation $\hat{d}$ and state estimation $\hat{x}(0)$. Write down the delta formulation that tracks $x_{\text{ref}}$ while accounting for such a disturbance. Will $\hat{d}$ appear in your problem formulation?\vspace{5\baselineskip}
    \item Transform the delta formulation in problem \ref{ex6:1} into a general QP form
    \begin{align*}
        \min_{\theta} \quad \theta^{\top}\Sigma\theta,\\
        \text{s.t.} \quad E\theta \leq e,
    \end{align*}
    which is better suited for the numerical implementation in optimization solvers. We write $\theta$ to be $[\Delta x_0^{\top},\cdots,\Delta x_N^{\top},\Delta u_0^{\top},\cdots,\Delta u_{N-1}^{\top}]^{\top}$. Assume $x_k \in \mathbb{R}^{n_x},u_k \in \mathbb{R}^{n_u}$.
    Write down the matrix $\Sigma$, $E$, and $e$, specify their dimensions and components for each block.
\end{enumerate}